\begin{lemma}[Scheide symmetric missing-color internal edges (13)--(14)]
Fix defective $\delta$.  Every color missing at one of the replacement
vertices in $U_\delta$ has an edge identity internal to
$U_{\gamma(\delta)}$; symmetrically, every color missing at one of the two
vertices of $U_{\gamma(\delta)}$ has an edge identity internal to
$U_\delta$.
\end{lemma}
\begin{proof}
Fix $\delta$, and abbreviate
\[
 X=\noderef{TreeSequenceVertexSet}(T),\qquad
 \gamma=\gamma(\delta),\qquad
 U_\gamma=\{u^1_\gamma,u^2_\gamma\},\qquad
 Z=X\setminus U_\gamma,
\]
and $U_\delta=\{u^1_\delta,u^2_\delta\}$.  Put
\[
 \varphi_\delta=
 \operatorname{swap}_{\gamma,\delta}
   (\varphi;v^0_\delta),\qquad
 X_\delta=V(T_\delta).
\]
Here $v^0_\delta,f^0_\delta,\alpha_0$, and the index defining $\gamma$
are the data of \noderef{ScheideDefectiveSetup}; the two additional
$\delta$-edges $f^1_\delta,f^2_\delta$, their outside ends
$u^1_\delta,u^2_\delta$, the replacement sequence $T_\delta$, and the
identity
\[
 X_\delta=Z\cup U_\delta                                      \tag{15}
\]
are the data of \noderef{ScheideReplacementData}.  Unpacking the extension
\noderef{ScheideReducedConfiguration}, the exact ordered layout and
attachment fields in \noderef{ScheideStructuralConfiguration}
show that the base $Y$ is contained in $Z$, that each of the two
$\gamma$-block edges joins $Y$ to the corresponding vertex of $U_\gamma$,
and that both are colored $\gamma$.  Since $v^0_\delta$ is a recorded exit,
\noderef{ScheideStructuralExitsInBase} also gives $v^0_\delta\in Y\subseteq
Z$.  In particular $y_0\in Y\subseteq Z$.

We first record all maximum-tree facts that will be used below.  The
inherited structural configuration makes $(T,e_0,\varphi)$ globally
maximum.  Applying \noderef{ScheideGlobalMaximumOrder} to this original
globally maximum triple gives
\[
 X\text{ is elementary and closed with respect to }\varphi.  \tag{16}
\]
The coloring $\varphi_\delta$ is exactly the component swap of
\noderef{TwoColorComponentSwap}.  By
\noderef{ScheideDefectiveSwitchEffect}, it is proper, leaves exactly
$e_0$ uncolored, changes the missing set at $v^0_\delta$ by deleting
$\gamma$ and inserting $\delta$, and fixes every other tree-incident edge
identity and every other old tree-vertex missing set.  Moreover
\noderef{ScheideReplacementMaximum} says that
$(T_\delta,e_0,\varphi_\delta)$ is globally maximum.  Applying
\noderef{ScheideGlobalMaximumOrder} to that triple, rather than treating
global maximality itself as an elementarity assertion, gives
\[
 X_\delta\text{ is elementary and closed with respect to }
 \varphi_\delta.                                               \tag{17}
\]

The colors $\gamma$ and $\delta$ are distinct.  Indeed, the setup says
that $\gamma$ is missing at $v^0_\delta$ while the incident edge
$f^0_\delta$ has color $\delta$.  Also $\alpha_0\ne\gamma$ by the
choice of $\alpha_0$ in \noderef{ScheideDefectiveSetup}.  Finally
$\alpha_0\ne\delta$: since $\delta$ is defective in the sense of
\noderef{DefectiveColors}, there is a $\delta$-colored boundary edge of
$X$; closedness in (16) therefore makes $\delta$ present at every vertex
of $X$, whereas $\alpha_0$ is missing at $y_0$.  It follows from the
exact missing-set formula in \noderef{ScheideDefectiveSwitchEffect} that
$\alpha_0$ remains missing at $y_0$ under $\varphi_\delta$: if
$y_0\ne v^0_\delta$ its entire missing set is fixed, and if
$y_0=v^0_\delta$ the formula deletes only $\gamma$ and inserts only
$\delta$.

We prove (13).  Let $\alpha$ be missing at one of
$u^1_\delta,u^2_\delta$ under $\varphi$.  The edge $f^i_\delta$ incident
with the selected vertex has color $\delta$, by
\noderef{ScheideReplacementData}, so $\alpha\ne\delta$.
The second edge supplied by
\noderef{ScheideSymmetricDefectiveInternal} is a $\gamma$-colored edge
whose ends are exactly $u^1_\delta,u^2_\delta$; hence $\gamma$ is present
at both vertices and $\alpha\ne\gamma$.  The agreement theorem
\noderef{ScheideReplacementPairFixed} says that $\varphi$ and
$\varphi_\delta$ agree on every edge incident with either replacement
vertex.  Thus $\alpha$ is missing at the same selected vertex under
$\varphi_\delta$.

That vertex and $y_0$ are distinct members of $X_\delta$: membership
follows from (15), and distinctness follows because the replacement data
put $U_\delta$ outside $X$ while $y_0\in X$.  Since $\alpha$ is missing
at the former and $\alpha_0$ is missing at the latter under
$\varphi_\delta$, elementariness in (17) gives
\[
 \alpha\ne\alpha_0.                                          \tag{18}
\]
Here and below ``elementary'' is used exactly in the sense of
\noderef{ColoredVertexSetElementary}: two distinct vertices of the set
cannot have a common missing color.
We have now proved separately that $\alpha$ differs from
$\alpha_0,\delta$, and $\gamma$.

Because $\alpha_0$ and $\alpha$ both differ from $\gamma$ and $\delta$,
the defining formula in \noderef{TwoColorComponentSwap} gives, for every
edge identity $q$,
\[
 \varphi(q)=\alpha_0\Longleftrightarrow
 \varphi_\delta(q)=\alpha_0,
 \qquad
 \varphi(q)=\alpha\Longleftrightarrow
 \varphi_\delta(q)=\alpha.                                  \tag{19}
\]
Thus the $(\alpha_0,\alpha)$ incidence graphs, and hence all their
components in the sense of \noderef{TwoColorComponent}, are identical
under $\varphi$ and $\varphi_\delta$.  Define $K_\gamma$ to be this
component under $\varphi$ through $u^1_\gamma$.

Apply \noderef{ScheideStructuralInternalPair} with the color $\alpha_0$.
Its hypotheses hold because $\alpha_0\ne\gamma$ and $\alpha_0$ is missing
at $y_0\in Z$.  It gives an $\alpha_0$-colored edge identity with ends
$u^1_\gamma,u^2_\gamma$.  Consequently both vertices of $U_\gamma$
belong to $K_\gamma$.  On the other hand, $K_\gamma$ contains no vertex
of $X_\delta$.  Indeed, $\alpha_0$ and $\alpha$ are missing at vertices
of the closed set $X_\delta$ by (17), so the definition
\noderef{ColoredVertexSetClosed} forbids a boundary edge of either color.
The component begins at $u^1_\gamma\notin X_\delta$ by (15); if it first
entered $X_\delta$, the entering edge of its incidence walk would be just
such a forbidden boundary edge.  Therefore
\[
 K_\gamma\cap X_\delta=\varnothing.                           \tag{20}
\]
In particular it avoids $Z\subseteq X_\delta$.  Every old tree edge has
an endpoint in $Z$: this is immediate for a base edge, and the exact block
attachments in \noderef{ScheideStructuralConfiguration} give a base end
in $Y\subseteq Z$ for every block edge.  Hence no edge identity of $T$
belongs to $K_\gamma$, for an incident edge-node in the component would
bring its $Z$-endpoint into the same component.

Now define explicitly
\[
 \psi=\operatorname{swap}_{\alpha_0,\alpha}
       (\varphi;u^1_\gamma)
\]
on the component $K_\gamma$.  By (18) and
\noderef{TwoColorComponentSwapProper}, $\psi$ is proper.  We next verify
the other parts of the original certificate separately.  The unique
uncolored identity remains $e_0$: the operation
\noderef{TwoColorComponentSwap} only interchanges two colored values and
fixes every other value, including the uncolored value.  Every old tree
edge keeps its color because no tree-edge identity belongs to
$K_\gamma$.

Every earlier missing-color witness needed by the Tashkinov condition also
survives.  For a nonfirst base edge, its earlier witness from the
Tashkinov-tree clause lies in the earlier part of the base $Y$.  For a
block edge, the structural configuration explicitly supplies a missing
witness in $Y$.  Since $Y\subseteq Z$ and (20) says that the component
avoids $Z$, no edge incident with any such witness lies in the component;
its full missing-color set is therefore unchanged.  These cases exhaust
the edges of the exact ordered layout, so $T$ remains a Tashkinov tree
under $\psi$ in the sense of \noderef{TashkinovTree}.

Unpack \noderef{GlobalMaximumTriple} and \noderef{KTriple}.  The graph,
the criticality and chromatic-index clauses, and the inequality involving
$k$ are unchanged; properness and the unique uncolored edge were just
checked; the preceding paragraph supplies the Tashkinov certificate; and
the vertex set and its cardinality are unchanged.  Thus
$(T,e_0,\psi)$ is again globally maximum.  In particular its cardinality
is $\noderef{TashkinovOrder}(G)$.

We show that no $\alpha$-edge can leave $U_\gamma$.  Any
$\alpha$-edge incident with either vertex of $U_\gamma$ belongs to
$K_\gamma$, because both of those vertices belong to the component.  Let
$q$ be such an edge and let $w$ be its other endpoint, whose existence
and distinction from the chosen endpoint follow from
\noderef{Incident} and the loopless law in
\noderef{LooplessMultigraph}.  If $w\in Z$, then $q$ is a boundary edge
of $X_\delta=Z\cup U_\delta$ in the sense of
\noderef{BoundaryEdges}.  By (19) it has color $\alpha$ under
$\varphi_\delta$, contradicting replacement closedness in (17), since
$\alpha$ is missing in $U_\delta$.

If instead $w\notin X$, then the switch changes $q$ from $\alpha$ to
$\alpha_0$.  Append the ordered pair $q,w$ to $T$.  The new vertex is
outside $X$, so it is distinct from every old vertex.  The new edge is
distinct from every tree edge because it has one end outside $X$, whereas
every tree edge has both ends in $X$ by \noderef{TreeSequence}; its other
end in $U_\gamma$ is an earlier attachment vertex.  Thus the appended
sequence is a tree-sequence.  All old Tashkinov conditions were preserved,
and the new $\alpha_0$-edge has the earlier witness $y_0$, whose missing
set was unchanged because $y_0\in Z$ is avoided by $K_\gamma$.  Hence the
appended sequence is a Tashkinov tree under $\psi$, with one more vertex
than $T$.  This contradicts the defining maximum in
\noderef{TashkinovOrder} and the globally maximum certificate just rebuilt.

The decomposition $X=Z\cup U_\gamma$ now shows that an $\alpha$-edge
incident with $U_\gamma$ either is internal to the pair or falls into one
of the two excluded cases.  If there were no internal $\alpha$-edge, both
vertices of $U_\gamma$ would therefore have no incident $\alpha$-edge,
so $\alpha$ would be missing at both by \noderef{MissingColors}.  They are
distinct vertices of $X$ by the structural layout and
\noderef{TreeSequence}, contradicting original elementariness in (16).
Thus an individual $\alpha$-colored edge has ends exactly
$u^1_\gamma,u^2_\gamma$, proving (13).

We now prove (14) without appealing merely to symmetry.  Let $\alpha$ be
missing at one of $u^1_\gamma,u^2_\gamma$ under $\varphi$.  The
corresponding block edge from
\noderef{ScheideStructuralConfiguration} is incident there and has color
$\gamma$, so $\alpha\ne\gamma$.  The first edge supplied by
\noderef{ScheideSymmetricDefectiveInternal} is a $\delta$-colored edge
with ends $u^1_\gamma,u^2_\gamma$, so $\alpha\ne\delta$.  Finally $y_0$
is distinct from both block vertices by the exact structural positions and
vertex injectivity in \noderef{TreeSequence}; since $\alpha_0$ is missing
at $y_0$ and $\alpha$ is missing at the selected block vertex,
elementariness in (16) gives $\alpha\ne\alpha_0$.  Thus in this direction
also $\alpha$ differs from all three of
$\alpha_0,\delta,\gamma$, and (19) holds for this $\alpha$ as well.

Apply \noderef{ScheideReplacementInternalMissing} to $\alpha_0$.  The
required witness is $y_0\in Z$, and $\alpha_0\ne\gamma$ was proved above.
We obtain an edge identity of original color $\alpha_0$ with ends
$u^1_\delta,u^2_\delta$.  By
\noderef{ScheideReplacementPairFixed}, this same edge has color
$\alpha_0$ under $\varphi_\delta$.  Define $K_\delta$ to be the
$(\alpha_0,\alpha)$ component under $\varphi_\delta$ through
$u^1_\delta$.  The internal $\alpha_0$-edge puts both vertices of
$U_\delta$ in $K_\delta$.

The component $K_\delta$ contains no vertex of the original set $X$.
Indeed, under $\varphi$, $\alpha_0$ is missing at $y_0\in X$ and
$\alpha$ is missing at a vertex of $U_\gamma\subseteq X$.  Original
closedness in (16) therefore forbids boundary edges of both colors.
By (19) the two-color incidence graph under $\varphi_\delta$ is the same
as under $\varphi$.  Since $U_\delta$ lies outside $X$, a first entry of
$K_\delta$ into $X$ would contradict that closure.  Consequently
\[
 K_\delta\cap X=\varnothing,
 \qquad K_\delta\cap X_\delta\subseteq U_\delta.               \tag{21}
\]

No edge identity of $T_\delta$ lies in $K_\delta$.  By the exact edge and
vertex records of \noderef{ScheideReplacementData}, each retained old
edge has an endpoint in $Z\subseteq X$, so (21) excludes its edge-node.
The only other tree edges are $f^1_\delta,f^2_\delta$.  They have original
color $\delta$ by the replacement data and the same color under
$\varphi_\delta$ by \noderef{ScheideReplacementPairFixed}; since
$\delta$ differs from $\alpha_0$ and $\alpha$, neither can be an edge-node
of this two-color component.

Define the second switched coloring explicitly by
\[
 \theta=\operatorname{swap}_{\alpha_0,\alpha}
       (\varphi_\delta;u^1_\delta)
\]
on $K_\delta$.  By \noderef{TwoColorComponentSwapProper}, applied now to
the proper coloring supplied by \noderef{ScheideReplacementMaximum},
$\theta$ is proper.  The unique uncolored edge remains $e_0$, because the
swap only interchanges the two colored values.  Every edge of
$T_\delta$ keeps its color because no such edge belongs to $K_\delta$.

We check the earlier missing witnesses for $T_\delta$ in two separate
classes.  A retained base edge keeps its original earlier base witness;
a retained block edge keeps the base witness given by
\noderef{ScheideStructuralConfiguration}.  Under the first switch from
$\varphi$ to $\varphi_\delta$, such a retained witness is not destroyed:
a retained base-edge color is different from every block color, and a
retained block color is different from the deleted color $\gamma$ by
injectivity of the block-color indexing.  This is also immediate from the
missing-set clauses of \noderef{ScheideDefectiveSwitchEffect}, separating
the case that the witness is $v^0_\delta$ from the case that it is not.
All these witnesses lie in $Y\subseteq Z$, so (21) says their missing sets
are unchanged by the second switch to $\theta$.

For the two appended edges, the check is different.  Both have color
$\delta$ under $\varphi_\delta$ by
\noderef{ScheideReplacementPairFixed}, and the exact switch formula in
\noderef{ScheideDefectiveSwitchEffect} says that $\delta$ is missing at
$v^0_\delta$.  This vertex lies in $Y\subseteq Z$ and is avoided by
$K_\delta$, so it remains a missing witness for both appended edges under
$\theta$.  Thus every edge position of the replacement sequence retains
an earlier missing-color witness, and $T_\delta$ is a Tashkinov tree under
$\theta$.

Unpacking the globally maximum replacement triple from
\noderef{ScheideReplacementMaximum} and then
\noderef{GlobalMaximumTriple}, the graph-theoretic and numerical fields are
unchanged; properness, the unique uncolored edge, and the Tashkinov
condition were just checked; and $T_\delta$ has the same maximum cardinality.
Therefore $(T_\delta,e_0,\theta)$ is globally maximum.

It remains to exclude edges leaving $U_\delta$.  Every
$\alpha$-edge under $\varphi_\delta$ incident with either replacement
vertex belongs to $K_\delta$, since the internal $\alpha_0$-edge put both
vertices in that component.  Let $q$ be such an edge with other endpoint
$w$.  If $w\in Z$, then
\noderef{ScheideReplacementPairFixed} gives $\varphi(q)=\alpha$.
The edge $q$ is a boundary edge of the original set $X$, because
$w\in Z\subseteq X$ and $U_\delta\cap X=\varnothing$.  This contradicts
original closedness in (16), since $\alpha$ is missing in
$U_\gamma\subseteq X$.

If $w\notin X_\delta$, then the switch from $\varphi_\delta$ to $\theta$
changes $q$ to color $\alpha_0$.  Append $q,w$ to $T_\delta$.  The same
field-by-field check as above applies: $w$ is a new vertex, $q$ is a new
edge because every old replacement-tree edge has both ends in
$X_\delta$, its replacement-pair endpoint is an earlier attachment
vertex, all old Tashkinov witnesses were preserved, and the appended edge
uses $y_0\in Z$ as its earlier missing $\alpha_0$ witness.  This produces
a Tashkinov tree under $\theta$ with one more vertex than $T_\delta$,
contradicting the globally maximum replacement certificate just rebuilt
and \noderef{TashkinovOrder}.

By (15), an $\alpha$-edge incident with $U_\delta$ is therefore either
internal to the pair or impossible.  If no internal one existed under
$\varphi_\delta$, both distinct vertices of $U_\delta$ would miss
$\alpha$, contradicting replacement elementariness in (17).  Hence there
is an edge $q$ with
\[
 \varphi_\delta(q)=\alpha,qquad
 G.\operatorname{ends}(q)=\{u^1_\delta,u^2_\delta\}.
\]
It is incident with both replacement vertices, so
\noderef{ScheideReplacementPairFixed} transfers its color back to
$\varphi(q)=\alpha$.  This is the required individual edge identity and
proves (14).
\end{proof}
